\section{Hybrid Experiments and Lemmas}
\label{sec:hybrids-details}

We compare the entire executions on the two sides of
Figure~\ref{fig:model-simulation}.  Every experiment runs an efficient
environment $\mathcal Z$ to its output bit, except at an explicitly
specified abort.  This includes its application
inputs, clock calls, ledger commands, corruptions, and local evaluations
of $\Sample$.  Application inputs are addressed to their pids; ledger
commands use the adversarial interface shown in the figure.  A reduction runs this whole experiment, replacing one
operation by a call to its primitive challenger.

The named replacements in Figure~\ref{fig:simulator-patches} give the
following route from the simulator to the protocol.  Each row changes
only the listed replacements; all other rules are inherited.
\begin{center}
\small
\renewcommand{\arraystretch}{1.15}
\begin{tabular}{@{}l>{\raggedright\arraybackslash}p{.47\linewidth}>{\raggedright\arraybackslash}p{.34\linewidth}@{}}
Experiment & Replacement & Property used \\\hline
$H_0$ & All replacements have subscript $0$ & Simulated execution \\
$H_1$ & $\InputEncryption_0\to\InputEncryption_1$ & Full NCE with tape exposure \\
$H_2$ & $\StateCiphertext_0\to\StateCiphertext_1$ & State-encryption IND-CPA \\
$H_3$ & $\OutputPads_0\to\OutputPads_1$;\newline $\StateCiphertext_1\to\StateCiphertext_2$ & Deferred decisions; correctness \\
$H_4$ & $\Garbling_0\to\Garbling_1$ & Adaptive garbling privacy \\
$H_5$ & $\Sharing_0\to\Sharing_1$ & Perfect sharing privacy and completion \\
$H_6$ & $\Deliveries_0\to\Deliveries_1$ & Receiver NCE with tape exposure \\
$H_7$ & Remove guards and ideal bookkeeping & Correctness and consistency \\
$H_8$ & $\Sampling_0\to\Sampling_1$ & Adaptive universal sampling
\end{tabular}
\end{center}

\subsection{Complete experiments}
\label{sec:hybrid-experiments}

Write $\SimExec^{\mathcal Z}$ for the ideal execution with the simulator
of Section~\ref{sec:simulator-details}, and $\RealExec^{\mathcal Z}$ for the real
execution with the dummy adversary.  Both use the clock of
Figure~\ref{fig:model-simulation}.  Let
\[
 p_i^{\mathcal Z}(\lambda)=\Pr[H_i^{\mathcal Z}(1^\lambda)=1].
\]
Probabilities include all coins of the experiment.  We do not condition
on committees satisfying their thresholds.

The first experiment $H_0$ runs $\SimExec$, but terminates with bit $0$
if one of the simulator's following abort conditions occurs:
\begin{description}
  \item[$\BadThreshold$.] An assigned committee has more than $t$
    initially corrupted members.
  \item[$\BadEarlyQuery$.] The eventual description $\SampleCircuit_s$
    has already been answered before the simulator reserves it.
  \item[$\BadPublicData$.] A consistency check in
    Section~\ref{sec:simulator-checks} fails.
  \item[$\BadAuth$.] A record is accepted before its provider's
    signing-key exposure whose message was never signed by the honest
    signing step.  Different signature bytes on an already signed
    message do not trigger this condition.
\end{description}
The intermediate experiments inherit these guards.  Thus $H_0$ is a
guarded version of the simulated execution; their output probabilities
can differ only when a guard is reached.  A further correctness guard
is introduced explicitly in $H_3$.  $H_8$ is the unguarded real execution.

Until the final step, execute the algorithms $\SimSetup$ and $\SimRO$
with the current memoized source $\SimSample$.  This defines the oracle
answers operationally even when the source returns modified samples.
We do not assume universal-sampler consistency for such a source;
Definition~\ref{def:adaptive-universal-sampler} supplies it once the
source has become honest.

Experiments $H_0$ through $H_6$ run the ideal task internally.  A
reduction can therefore read the secret inputs, private outputs, and
states $\sigma_s$ needed to define a replacement.  These values are not
made available to $\mathcal Z$ or to the implementable simulator.
The circuit rule remains
$F_s=G(s,(\xpub_t,\ypub_t)_{t\leq s-2})$, using the committed public
input components.  Announcements, commitments, and private deliveries
retain their specified clock times, including queued outputs delivered
to a pid corrupted before the delivery tick.

\paragraph{Input authentication and fresh descriptions.}
The authentication guard is bounded using one-time signature
unforgeability.  Its reduction stops at the first new accepted message
under its unexposed challenge signer and need not answer a later
corruption.  An honest record already posted cannot be replaced by a
later record, even after the signing key is exposed.

For nonce freshness, consider the execution stopped before the first
honest nonce-committee key is disclosed.  Encryption privacy hides the
honest share packets, and threshold privacy hides $Q_s$ given the
corrupt shares.  The description's other fields are frozen before
release, so an earlier request for that description must predict its
nonce.  Count every query to the sample source, including those made
by $\SimSetup$ and $\SimRO$.  This argument uses receiver NCE to hide
the packets only in the stopped execution; it does not have to explain
them after disclosure.  The guard bounds, including threshold failure,
are separate obligations in Lemma~\ref{lem:simulation-guards} below.

\paragraph{Executable replacements.}
$H_0$ uses the seven subscript-$0$ hooks of
Section~\ref{sec:simulator-hooks}.  Each following experiment inherits
all logistics, checks, and hook versions from its predecessor, except
for the replacements specified below.  These definitions use local
algorithms; challenger commands are used only in reductions.

\begin{description}
\item[$H_1$: $\InputEncryption_1$.]
 $\Prepare(s,a)$ generates a genuine $\SNCE$ key pair and retains its
 key-generation tape.  It samples the same independent decryption tape
 $\delta_{s,a}$ and returns $\inputpk_{s,a}$.
 In $\Post(s,a,\xpub)$, replace the reserved ciphertext by
 \[
 c_{s,a}\gets\SNCE.\Enc(\inputpk_{s,a},\xsec_{s,a};\rho_{s,a}),
 \]
 where $\xsec_{s,a}$ is the selected secret input of the internally run
 ideal task and $\rho_{s,a}$ is freshly sampled and retained.  Keep the
 record construction, signing, and honest-posting flag updates.
 $\Corrupt(s,a,m)$ uses the retained sender tape after an honest posting,
 and the genuine receiver key otherwise; it makes no NCE explanation.
 $\Extract$ is unchanged.  The receiver key and its generation tape
 remain internal, not part of the input provider's disclosure.

\item[$H_2$: $\StateCiphertext_1$.]
 The call $\StateCiphertext_1(s)$ samples fresh coins $\rho_s$ and returns
 \[
 c_s\gets\PKE.\Enc(\statepk_{s+1},\sigma_s;\rho_s),
 \]
 using the state computed by the internal ideal task.  It replaces the
 zero-state ciphertext wherever the input-closing logistics uses it.

\item[$H_3$: $\OutputPads_1$ and $\StateCiphertext_2$.]
 During sample preparation, before the garbling hook, sample and retain
 all pads $r_{s,j}$ and the state-encryption tape $\rho_s$.
 $\OutputPads_1.\Prepare$ passes each actual pad packet to
 $\Deliveries.\Packet$, including for honest recipients.
 $\Close$ sets $z_s=\ysec_s\oplus r_s$ using the internally computed
 ideal output.  Every later pad binding uses the retained pad.
 $\StateCiphertext_2(s)$ returns the preceding encryption using the
 retained $\rho_s$, rather than sampling new coins.

 Before the encoding call at input closure, abort with bit $0$ if a
 selected honest ciphertext does not decrypt under its retained input
 key and $\delta_{s,a}$ to its committed secret input, or if
 $\PKE.\Dec(\statesk_{s+1},c_s)\ne\sigma_s$.
 Call this guard $\BadCorrectness$.  Corrupt ciphertexts are not subject
 to the first test: their ideal inputs were defined by $\Extract$.

\item[$H_4$: $\Garbling_1$.]
 $\Prepare(s)$ constructs
 \[
 C_s=\RoundCircuit_s[\statesk_s,\boldsymbol{\inputsk}_s,
                    \boldsymbol\delta_s,\boldsymbol r_s,\rho_s]
 \]
 with the already retained constants, runs
 $(\GC_s,e_s)\sample\Garble(1^\lambda,C_s)$, and stores $e_s$.
 $\Encode(s,(z_s,c_s,\ypub_s))$ instead computes
 $X_s\gets\Encode(e_s,\InputBits_s)$ from the public input string fixed
 by the logistics.  Retain the same prescribed result for the existing
 public-evaluation check.  No garbling-simulator algorithm is called.

\item[$H_5$: $\Sharing_1$.]
 In $\Prepare(s)$, retain the original uniform shares at initially
 corrupt positions.  Extract both labels and the common material from
 the projective encoding key $e_s$.  For every $(w,b)$ independently
 complete a full sharing to
 $\langle\GarbleAux_s,L_{s,w,b}\rangle$, conditioned on those shares.
 Cache both vectors and pass every member's actual encoded share to
 $\Deliveries.\Packet$, replacing the $\bot$ arguments for honest
 members.  $\Open$ uses the cached selected vector and binds its honest
 members' deliveries without completing the sharing again.

\item[$H_6$: $\Deliveries_1$.]
 $\Register$ uses genuine $\RNCE.\KeyGen$ and retains its complete tape
 for every delivery key.  For application registrations, change the
 internally run ideal registry's key-generation step too: use this
 same key and the same genuine signing pair in both registries.
 The internal registration reply now supplies its real receiver key
 and tape, together with the signing material, to $\Deliveries_1.\Register$.
 No NCE helper is generated for this key.
 $\Packet(P,m)$ now receives a known plaintext at every designated call;
 it samples and caches $\RNCE.\Enc(\ncek_P,m)$.
 $\Bind(P,m)$ simply returns the retained receiver key.
 Histories use the retained key-generation tape; there are no receiver
 explanations.  Nonce packets and other sample descriptions continue
 to use ordinary encryption.

\item[$H_7$: remove the ideal bookkeeping.]
 Run the genuine $\FLedger[\Pi]$ with the dummy adversary and
 $\mathcal Z$, but obtain public parameters and oracle answers from
 $\SimSetup^{\IdealSample}$ and $\SimRO^{\IdealSample}$.
 Here $\IdealSample(d)$ draws and memoizes an honest $d(\omega)$.
 There is no internal ideal task, source-tape access by the surrounding
 execution, or hybrid guard.  The comparison from $H_6$ uses correctness
 to replace the ideal bookkeeping by the protocol's own outputs and
 retained histories.

\item[$H_8$: $\Sampling_1$.]
 Replace the simulated public parameters and oracle in $H_7$ by genuine
 setup and the random oracle.  This is $\RealExec$.
 The surrounding execution uses only public sample values and registered
 party state; it does not request sampling tapes or sample-owned secret
 keys from the universal-sampler challenger.
\end{description}

\subsection{Lemma sequence}
\label{sec:hybrid-lemmas}

All statements below quantify over the same complete experiments and
arbitrary PPT environments $\mathcal Z$.  Each negligible function may
depend on $\mathcal Z$.  We give the full reduction for
Lemma~\ref{lem:restore-input-encryption} in
Section~\ref{sec:first-hybrid-proof}.  The remaining statements specify
the proof obligations; their reductions are still to be written.
Section~\ref{sec:hybrids} records the proposed arguments.


\begin{lemma}[Restoring input encryption]
\label{lem:restore-input-encryption}
Assume full non-committing security with tape exposure
(Definition~\ref{def:nce}).  For every PPT $\mathcal Z$,
\[
 |p_0^{\mathcal Z}(\lambda)-p_1^{\mathcal Z}(\lambda)|
 \leq\negl(\lambda).
\]
\end{lemma}

\begin{lemma}[Restoring the encrypted states]
\label{lem:restore-state-ciphertexts}
Assume IND-CPA security of $\PKE$ (Definition~\ref{def:ind-cpa}).
For every PPT $\mathcal Z$,
\[
 |p_1^{\mathcal Z}(\lambda)-p_2^{\mathcal Z}(\lambda)|
 \leq\negl(\lambda).
\]
\end{lemma}

\begin{lemma}[Fixing the circuit constants]
\label{lem:fix-circuit-constants}
Let $H_3^-$ make the early-pad and early-coin replacements defining
$H_3$, without its new correctness guard.  For every PPT $\mathcal Z$,
\[
 p_2^{\mathcal Z}(\lambda)
   =\Pr[(H_3^-)^{\mathcal Z}(1^\lambda)=1].
\]
Under correctness of input and state encryption,
$|p_2^{\mathcal Z}(\lambda)-p_3^{\mathcal Z}(\lambda)|
\leq\negl(\lambda)$.
\end{lemma}

\begin{lemma}[Restoring the garblings]
\label{lem:restore-garblings}
Assume adaptive garbling privacy
(Definition~\ref{def:adaptive-garbling}).  For every PPT $\mathcal Z$,
\[
 |p_3^{\mathcal Z}(\lambda)-p_4^{\mathcal Z}(\lambda)|
 \leq\negl(\lambda).
\]
\end{lemma}

\begin{lemma}[Preparing the committee sharings]
\label{lem:restore-sharings}
Under perfect sharing privacy and the conditional completion property
of Lemma~\ref{lem:sharing-completion}, for every PPT $\mathcal Z$,
\[
 p_4^{\mathcal Z}(\lambda)=p_5^{\mathcal Z}(\lambda).
\]
\end{lemma}

\begin{lemma}[Restoring private deliveries]
\label{lem:restore-deliveries}
Assume receiver non-committing security with tape exposure
(Definition~\ref{def:receiver-nce}).  For every PPT $\mathcal Z$,
\[
 |p_5^{\mathcal Z}(\lambda)-p_6^{\mathcal Z}(\lambda)|
 \leq\negl(\lambda).
\]
\end{lemma}

\begin{lemma}[Removing the ideal bookkeeping]
\label{lem:remove-ideal-bookkeeping}
Assume correctness of the cryptographic primitives, robust reconstruction,
adaptive universal-sampler consistency, and the assumptions of
Lemma~\ref{lem:simulation-guards}.  For every PPT $\mathcal Z$,
\[
 |p_6^{\mathcal Z}(\lambda)-p_7^{\mathcal Z}(\lambda)|
 \leq\negl(\lambda).
\]
This comparison removes the guards and replaces the internal ideal task
by the protocol's own application outputs and retained histories.
\end{lemma}

\begin{lemma}[Restoring public universal sampling]
\label{lem:restore-public-sampling}
Assume adaptive universal-sampler security.  For every PPT $\mathcal Z$,
\[
 |p_7^{\mathcal Z}(\lambda)-p_8^{\mathcal Z}(\lambda)|
 \leq\negl(\lambda),
 \qquad
 p_8^{\mathcal Z}(\lambda)=\Pr[\RealExec^{\mathcal Z}(1^\lambda)=1].
\]
\end{lemma}

\begin{lemma}[From the simulator to its guarded experiment]
\label{lem:simulation-guards}
Assume the primitive properties used in
Lemmas~\ref{lem:restore-input-encryption}--\ref{lem:restore-deliveries},
adaptive universal-sampler security, one-time signature security, and
correctness of the primitives.
Choose the committee parameters so that the
probability of an over-threshold committee in any permitted efficient
execution is negligible.  Then, for every PPT $\mathcal Z$,
\[
 \left|\Pr[\SimExec^{\mathcal Z}(1^\lambda)=1]
             -p_0^{\mathcal Z}(\lambda)\right|\leq\negl(\lambda).
\]
\end{lemma}

For the consistency guard, first bound failure in the honest-source
experiment $H_6$.  To transfer the bound to $H_0$, let $\mathcal Z^1$
run $\mathcal Z$ with the same interaction but return $1$ at termination.
Then
\[
 p_i^{\mathcal Z^1}(\lambda)
   =1-\Pr[\text{a guard fires before termination in }H_i^{\mathcal Z}],
 \qquad 0\leq i\leq6.
\]
The preceding whole-experiment comparisons apply to $\mathcal Z^1$
without assuming the early guards negligible.  This transfers the
honest-$H_6$ bound backwards and avoids assuming consistency of a
modified sample source.  Establishing the direct $H_6$ bound is part
of the remaining proof of Lemma~\ref{lem:simulation-guards}.

Once these lemmas are proved, the endpoint comparison is just
\[
\begin{split}
 &\left|\Pr[\SimExec^{\mathcal Z}(1^\lambda)=1]
       -\Pr[\RealExec^{\mathcal Z}(1^\lambda)=1]\right|\\
 &\quad\leq
 \left|\Pr[\SimExec^{\mathcal Z}(1^\lambda)=1]-p_0^{\mathcal Z}(\lambda)\right|
 +\sum_{i=0}^{7}|p_i^{\mathcal Z}(\lambda)-p_{i+1}^{\mathcal Z}(\lambda)|.
\end{split}
\]
Each cryptographic comparison sweeps over the polynomially many relevant
instances reached by the fixed efficient experiment.  Its reduction may
use a polynomial bound on that number; neither setup nor the simulator
receives such a bound.  Within a sweep, count instances when they first
occur and continue all subsequent adaptive choices normally.

\subsection{First reduction: input encryption}
\label{sec:first-hybrid-proof}
\begin{proof}[Lemma~\ref{lem:restore-input-encryption}]
Fix $\mathcal Z$.  Enumerate input-key instances when their designated
round samples are first prepared, using provider order within a sample.
A repeated request for the same description returns the stored sample
and creates no new instance.  Let $J_j$ restore genuine input encryption
for the first $j$ instances and use the simulated procedure for the
rest.  Thus $J_0=H_0$ and $J_N=H_1$, where $N=N(\lambda)\geq1$ is a polynomial
bound on the number of instances in this experiment.  Such a bound is
used only by the reduction, not by the simulator or protocol.

We construct one uniform NCE adversary $\mathcal B$, receiving the
public challenge length $\ell$.  It chooses $j$ uniformly from $[N]$
and runs the experiment containing $\mathcal Z$, the clock, the ideal
computation, and the simulated ledger.  It also runs $\SimSetup$,
$\SimRO$, and their memoized sampling source.  Before the $j$th input-key
instance, it follows the common prefix of $J_{j-1}$ and $J_j$ without
using $\ell$.  If that instance never occurs, or its prescribed length
is different from $\ell$, it calls $\Finish(0)$.  Otherwise it replaces the target calls to
$\InputEncryption.\Prepare$, $\Post$, and $\Corrupt$ by the
following challenger calls.

\begin{description}
  \item[Preparation.]
    Call $\GetPublicKey()$ and publish the returned key as
    $\inputpk_{s,a}$.  Independently sample the decryption tape
    $\delta_{s,a}$.  Prepare the simulated garbling and the remaining
    sample components as prescribed.  No input-decryption secret key
    is needed: $\SimGarble$ receives only the circuit topology, and
    $\SimEncode$ will receive only the prescribed evaluation result.

  \item[Honest posting first.]
    When provider $a$ posts its selected honest input, call
    \[
      c\gets\EncryptMessage(\xsec_{s,a}).
    \]
    Sign and post the returned ciphertext
    together with the unchanged public component $\xpub_{s,a}$, using
    the genuine signing key and fresh signing coins.  If the provider
    is subsequently corrupted, call $\CorruptSender()$ and use its
    returned encryption tape in the disclosed history.  Repeated
    disclosures reuse that tape.  The input-decryption key and its
    generation tape are not part of this provider's state.

  \item[Corruption first.]
    As soon as both preparation and corruption have occurred, call
    $\EncryptMessage(0^\ell)$, keep the ciphertext private, and call
    $\CorruptReceiver()$.  Retain the returned secret key internally
    and discard the challenge ciphertext.  Interpret adversarial
    postings with this key and $\delta_{s,a}$, applying the specified
    decryption-failure rule, and submit the resulting secret input and
    the record's public component to the internal ideal computation.
    This uses no additional encryption challenge.  The provider's
    corruption disclosure contains its registered state and retained
    history, not this internal decryption key.
\end{description}
If neither posting nor corruption occurs, no encryption call is
required.  The branches also cover a provider corrupt at preparation
and corruption after an activation at which the provider stayed
silent.  Every challenge call is legal: there is one public-key call,
at most one encryption call, and any corruption call follows that
encryption.  Only one branch is used, with the same full-NCE simulator
underlying both.  The reduction then continues the remaining experiment
and calls $\Finish$ with the terminal bit of $\mathcal Z$, or with
$0$ if an inherited guard fires.

Every inherited check is executable in this reduction.  In particular,
sample and public-evaluation comparisons use cached values and public
computations, while sender-history reconstruction uses the explanation
returned by $\CorruptSender$ and the separate registered keys.  None
requires the honest-posting branch's unrevealed input-decryption key.

The reduction knows $\xsec_{s,a}$ at an honest posting because it runs
the ideal computation and its input interface internally.  This is
knowledge of the reduction, not an additional interface for the
implementable simulator.  All information disclosed to $\mathcal Z$
continues to follow the hybrid's interfaces.  The reduction knows the
provider's independent signing and delivery state and can therefore
complete its corruption disclosure.

In the ideal NCE experiment the challenger already sampled the reserved
ciphertext at $\GetPublicKey$.  Leaving it there until
$\EncryptMessage$ is exact: before either branch, the hybrid uses
neither that ciphertext nor its NCE simulator state.  They do not enter
the public sample or the topology-only garbling.  In the corruption-first
branch, the challenger's cached dummy-message explanation gives exactly
the key used by $J_{j-1}$.  In the real experiment it gives the genuine
receiver key.  In the honest-posting branch the two challenge worlds
supply respectively the simulated and genuine ciphertext and sender
tape.  All other instances are generated locally in their prescribed
modes.  Ordinary sample descriptions are evaluated as $d(\rho)$, and
computations involving the challenge public key or public ciphertext
require no unavailable key or simulator state.  This includes
$\mathcal Z$'s local sampling and interleaved oracle queries.  Consequently,
on a matching target length the two challenge worlds reproduce the
corresponding branches of $J_{j-1}$ and $J_j$.

Let $b(\lambda)$ bound all input lengths in this experiment, and write
$\Delta_\ell$ for the signed real-minus-ideal NCE advantage of
$\mathcal B$ at length $\ell$.  Target length is determined before the
challenge call and its events partition the common prefix.  Hence
\[
\begin{aligned}
  \sum_{\ell=0}^{b(\lambda)}\Delta_\ell
    &=\frac1N\sum_{j=1}^{N}
       \bigl(\Pr[J_j^{\mathcal Z}(1^\lambda)=1]
             -\Pr[J_{j-1}^{\mathcal Z}(1^\lambda)=1]\bigr)\\
    &=\frac{p_1^{\mathcal Z}(\lambda)-p_0^{\mathcal Z}(\lambda)}N.
\end{aligned}
\]
Uniform full-NCE security bounds every $|\Delta_\ell|$ by the same
negligible $\varepsilon_{\mathcal B,b}$.  The desired difference is
therefore at most $N(b+1)\varepsilon_{\mathcal B,b}$, which is negligible.
\end{proof}
